% GENERATED FILE. Edit canonical lessons, then run npm run generate:books.
% canonical-source-hash: fnv1a64:8c4ebbbcb548f678
% canonical-lessons: MR-C534-giryarohan, MR-C534-jaltaran, MR-C534-sibir, MR-C534-tambu, MR-C534-sekoti

\chapter{Into the Outdoors}
\label{ch:mr-into-the-outdoors}
\hlchaptermodality{534}

\begin{chapteropening}
\textbf{By the end of this chapter:} \emph{I can talk about climbing, swimming and camping.}

Complete the last lesson of chapter 534: I can talk about climbing, swimming and camping.
\end{chapteropening}

\section[\texorpdfstring{\mr{गिर्यारोहण} (giryārohaṇ)}{गिर्यारोहण (giryārohaṇ)}]{\mr{गिर्यारोहण} (giryārohaṇ) — mountaineering}

\label{lesson:MR-C534-giryarohan}



\begin{quote}
Before the new one: say the Marathi for to bend, then the Marathi for to leap.
\end{quote}

\subsection*{You'll want to know: \mr{गिर्यारोहण}}
\textbf{\mr{गिर्यारोहण}} — \emph{giryārohaṇ} — \textquotedblleft{}mountaineering\textquotedblright{}.

\begin{quote}
\mr{माझ्या} \mr{भावाला} \mr{गिर्यारोहणाची} \mr{आवड} \mr{आहे}.
\end{quote}

\begin{quote}
\emph{mājhyā bhāvālā giryārohaṇācī āvaḍ āhe.}
\end{quote}

\begin{quote}
My brother loves mountaineering.
\end{quote}

Neuter.

\begin{cousinweb}
\textbf{\mr{गिर्यारोहण}} (\emph{giryārohaṇ}) is a tatsama compound of Sanskrit \emph{giri}, \textquotedblleft{}mountain\textquotedblright{}, and \emph{ārohaṇa}, \textquotedblleft{}climbing\textquotedblright{} (Monier-Williams).
\end{cousinweb}

\begin{grammarlens}[title={What you've built}]
One more word for describing what you lived and felt.
\end{grammarlens}

\subsection*{Your turn}
\begin{itemize}
\raggedright
  \item \emph{Say it:} \emph{giryārohaṇ}
  \item \emph{Say it:} \emph{giryārohaṇ}, once more
  \item \emph{Say it:} \emph{mājhyā bhāvālā giryārohaṇācī āvaḍ āhe.} — My brother loves mountaineering.
  \item \emph{Recall:} say the Marathi for to bend, then the Marathi for to leap, then say \emph{giryārohaṇ} again
\end{itemize}

\begin{tcolorbox}[breakable,colback=teal!4,colframe=teal!35!black,title={Before you move on}]
What does \mr{गिर्यारोहण} mean? (\textquotedblleft{}Mountaineering\textquotedblright{}.) Say it once more.
\end{tcolorbox}
\section[\texorpdfstring{\mr{जलतरण} (jaltaraṇ)}{जलतरण (jaltaraṇ)}]{\mr{जलतरण} (jaltaraṇ) — swimming as a sport}

\label{lesson:MR-C534-jaltaran}



\begin{quote}
Before the new one: say the Marathi for to leap, then the Marathi for mountaineering.
\end{quote}

\subsection*{You'll want to know: \mr{जलतरण}}
\textbf{\mr{जलतरण}} — \emph{jaltaraṇ} — \textquotedblleft{}swimming as a sport\textquotedblright{}.

\begin{quote}
\mr{जलतरणाच्या} \mr{स्पर्धेत} \mr{मीरा} \mr{पहिली} \mr{आली}.
\end{quote}

\begin{quote}
\emph{jaltaraṇācyā spardhet mīrā pahilī ālī.}
\end{quote}

\begin{quote}
Meera came first in the swimming competition.
\end{quote}

Neuter.

\begin{cousinweb}
\textbf{\mr{जलतरण}} (\emph{jaltaraṇ}) is a tatsama compound of Sanskrit \emph{jala}, \textquotedblleft{}water\textquotedblright{}, and \emph{taraṇa}, \textquotedblleft{}crossing\textquotedblright{} (Monier-Williams).
\end{cousinweb}

\begin{grammarlens}[title={What you've built}]
One more word for describing what you lived and felt.
\end{grammarlens}

\subsection*{Your turn}
\begin{itemize}
\raggedright
  \item \emph{Say it:} \emph{jaltaraṇ}
  \item \emph{Say it:} \emph{jaltaraṇ}, once more
  \item \emph{Say it:} \emph{jaltaraṇācyā spardhet mīrā pahilī ālī.} — Meera came first in the swimming competition.
  \item \emph{Recall:} say the Marathi for to leap, then the Marathi for mountaineering, then say \emph{jaltaraṇ} again
\end{itemize}

\begin{tcolorbox}[breakable,colback=teal!4,colframe=teal!35!black,title={Before you move on}]
What does \mr{जलतरण} mean? (\textquotedblleft{}Swimming as a sport\textquotedblright{}.) Say it once more.
\end{tcolorbox}
\section[\texorpdfstring{\mr{शिबिर} (śibir)}{शिबिर (śibir)}]{\mr{शिबिर} (śibir) — a camp}

\label{lesson:MR-C534-sibir}



\begin{quote}
Before the new one: say the Marathi for mountaineering, then the Marathi for swimming as a sport.
\end{quote}

\subsection*{You'll want to know: \mr{शिबिर}}
\textbf{\mr{शिबिर}} — \emph{śibir} — \textquotedblleft{}a camp\textquotedblright{}.

\begin{quote}
\mr{सुट्टीत} \mr{मुलांसाठी} \mr{जंगलात} \mr{शिबिर} \mr{होतं}.
\end{quote}

\begin{quote}
\emph{suṭṭīt mulāṁsāṭhī jaṅglāt śibir hotaṁ.}
\end{quote}

\begin{quote}
In the holidays there was a camp in the forest for children.
\end{quote}

Neuter.

\begin{cousinweb}
\textbf{\mr{शिबिर}} (\emph{śibir}) is a tatsama, taken from Sanskrit \emph{śibira}, \textquotedblleft{}a camp, a tent\textquotedblright{} (Monier-Williams).
\end{cousinweb}

\begin{grammarlens}[title={What you've built}]
One more word for describing what you lived and felt.
\end{grammarlens}

\subsection*{Your turn}
\begin{itemize}
\raggedright
  \item \emph{Say it:} \emph{śibir}
  \item \emph{Say it:} \emph{śibir}, once more
  \item \emph{Say it:} \emph{suṭṭīt mulāṁsāṭhī jaṅglāt śibir hotaṁ.} — In the holidays there was a camp in the forest for children.
  \item \emph{Recall:} say the Marathi for mountaineering, then the Marathi for swimming as a sport, then say \emph{śibir} again
\end{itemize}

\begin{tcolorbox}[breakable,colback=teal!4,colframe=teal!35!black,title={Before you move on}]
What does \mr{शिबिर} mean? (\textquotedblleft{}A camp\textquotedblright{}.) Say it once more.
\end{tcolorbox}
\section[\texorpdfstring{\mr{तंबू} (tambū)}{तंबू (tambū)}]{\mr{तंबू} (tambū) — a tent}

\label{lesson:MR-C534-tambu}



\begin{quote}
Before the new one: say the Marathi for swimming as a sport, then the Marathi for a camp.
\end{quote}

\subsection*{You'll want to know: \mr{तंबू}}
\textbf{\mr{तंबू}} — \emph{tambū} — \textquotedblleft{}a tent\textquotedblright{}.

\begin{quote}
\mr{शिबिरात} \mr{आम्ही} \mr{तंबूत} \mr{झोपलो}.
\end{quote}

\begin{quote}
\emph{śibirāt āmhī tambūt jhoplo.}
\end{quote}

\begin{quote}
At the camp we slept in a tent.
\end{quote}

Masculine.

\begin{cousinweb}
\textbf{\mr{तंबू}} (\emph{tambū}) is usually traced to Persian; the exact path into Marathi is not settled.
\end{cousinweb}

\begin{grammarlens}[title={What you've built}]
One more word for describing what you lived and felt.
\end{grammarlens}

\subsection*{Your turn}
\begin{itemize}
\raggedright
  \item \emph{Say it:} \emph{tambū}
  \item \emph{Say it:} \emph{tambū}, once more
  \item \emph{Say it:} \emph{śibirāt āmhī tambūt jhoplo.} — At the camp we slept in a tent.
  \item \emph{Recall:} say the Marathi for swimming as a sport, then the Marathi for a camp, then say \emph{tambū} again
\end{itemize}

\begin{tcolorbox}[breakable,colback=teal!4,colframe=teal!35!black,title={Before you move on}]
What does \mr{तंबू} mean? (\textquotedblleft{}A tent\textquotedblright{}.) Say it once more.
\end{tcolorbox}
\section[\texorpdfstring{\mr{शेकोटी} (śekoṭī)}{शेकोटी (śekoṭī)}]{\mr{शेकोटी} (śekoṭī) — a campfire; a bonfire}

\label{lesson:MR-C534-sekoti}



\begin{quote}
Before the new one: say the Marathi for a camp, then the Marathi for a tent.
\end{quote}

\subsection*{You'll want to know: \mr{शेकोटी}}
\textbf{\mr{शेकोटी}} — \emph{śekoṭī} — \textquotedblleft{}a campfire; a bonfire\textquotedblright{}.

\begin{quote}
\mr{रात्री} \mr{सगळे} \mr{शेकोटीभोवती} \mr{बसले}.
\end{quote}

\begin{quote}
\emph{rātrī sagḷe śekoṭībhovtī basle.}
\end{quote}

\begin{quote}
At night everyone sat around the campfire.
\end{quote}

Feminine.

\begin{cousinweb}
\textbf{\mr{शेकोटी}} (\emph{śekoṭī}) is built on \mr{शेकणे}, \textquotedblleft{}to warm oneself at a fire\textquotedblright{}: CDIAL lists Marathi \emph{śekṇẽ} under \emph{*sēkk}, \textquotedblleft{}heat, foment, roast\textquotedblright{} (CDIAL 13581). The ending -\mr{ओटी} is not explained there.
\end{cousinweb}

\begin{grammarlens}[title={What you've built}]
One more word for describing what you lived and felt.
\end{grammarlens}

\subsection*{Your turn}
\begin{itemize}
\raggedright
  \item \emph{Say it:} \emph{śekoṭī}
  \item \emph{Say it:} \emph{śekoṭī}, once more
  \item \emph{Say it:} \emph{rātrī sagḷe śekoṭībhovtī basle.} — At night everyone sat around the campfire.
  \item \emph{Recall:} say the Marathi for a camp, then the Marathi for a tent, then say \emph{śekoṭī} again
\end{itemize}

\begin{tcolorbox}[breakable,colback=teal!4,colframe=teal!35!black,title={Before you move on}]
What does \mr{शेकोटी} mean? (\textquotedblleft{}A campfire; a bonfire\textquotedblright{}.) Say it once more.
\end{tcolorbox}
